\subsection[Closures]{Closure with respect to the rest of the world: \GMR, \SOE, and \SATE}\label{app:closures}

This subsection is new relative to \citet{LanzRutherford2016}.
It documents how the macro \texttt{loadrm(rr)} (file \texttt{code/loadmdl.gms}) and the flags \texttt{sate} and \texttt{sate2} (file \texttt{code/gtap12data.gms}) select the trade closure, and states the equilibrium conditions which are added to or replace those of Appendix \ref{equil}.
Let $R^M$ (\texttt{rm(r)}) denote the regions that are part of the model and $R^X$ (\texttt{rx(r)}) the remaining regions which form the rest of the world (\textsc{row}).
The \GMR model is obtained with $R^M$ equal to all regions and $R^X=\emptyset$; all conditions of Appendix \ref{equil} then apply to all regions and no additional variables are active.

\subsubsection{Variables and parameters}

Table \ref{tab:closurevars} lists the variables that are active only in the small open economy closures.
In addition, the following benchmark parameters are computed by \texttt{loadrm} for $r\in R^M$: exports to the rest of the world, $\textit{vem}_{ir}=\sum_{s\in R^X}\textit{vxmd}_{irs}$ [\texttt{vem(i,r)}]; the trade-weighted average export subsidy rate on these exports, $t^{xs,row}_{ir}=\sum_{s\in R^X}\textit{vxmd}_{irs}\,t^{xs}_{irs}/\textit{vem}_{ir}$ [\texttt{rtxs\_row(i,r)}], and the associated reference price $\textit{pem}^0_{ir}=1-\overline{t}^{xs,row}_{ir}$ [\texttt{pem0(i,r)}].
For $s\in R^X$, $\textit{vem}_{is}=\sum_{r\in R^M}\textit{vxmd}_{isr}$ denotes the exports of \textsc{row} region $s$ to the modelled regions.

\begin{table}[t!]
\begin{center}
\caption{Variables of the small open economy closures.}\label{tab:closurevars}
\small
\begin{tabular}{lp{6.2cm}lr}
\toprule
Variable & Definition & \GAMS variable & Active in \\
\midrule
$X_{ir}$, $r\in R^M$ & Exports of model region $r$ to \textsc{row} & \texttt{X(i,r)} & \SOE, \SATE \\
$X_{is}$, $s\in R^X$ & Exports of \textsc{row} region $s$ to the model & \texttt{X(i,s)} & \SOE, \SATE \\
$p^{FX}$ & Price of foreign exchange (real exchange rate) & \texttt{PFX} & \SOE, \SATE \\
$\textit{ROW}$ & Income of the rest-of-world agent & \texttt{ROW} & \SOE, \SATE \\
\addlinespace
$p^{PX}_{ir}$ & Price of exports of $r$ in its export market & \texttt{PPX(i,r)} & \texttt{sate} \\
$E_{ir}$ & Export demand index (benchmark $=1$) & \texttt{EXPORTS(i,r)} & \texttt{sate} \\
$\textit{RX}_{ir}$ & Foreign exchange earned from exports (index) & \texttt{REVEXP(i,r)} & \texttt{sate} \\
\addlinespace
$p^{FIX}_{ir}$ & Price of the \textsc{row}-owned fixed factor in export demand & \texttt{PFIX(i,r)} & \texttt{sate2} \\
\bottomrule
\end{tabular}
\end{center}
\end{table}

\subsubsection{Small open economy with fixed world prices (\SOE)}

Exports of a modelled region to the rest of the world are converted into foreign exchange at a fixed rate.
With $p^X_{ir}$ denoting the export supply price (equal to $p^Y_{ir}$ if $\textit{etrndx}_i=\infty$), the zero profit condition of $X_{ir}$ reads
\begin{equation}\label{zpX}
p^X_{ir}\,\frac{1-t^{xs,row}_{ir}}{1-\overline{t}^{xs,row}_{ir}} \;\geq\; p^{FX}
\qquad \perp \; X_{ir}\geq 0, \quad r\in R^M,
\end{equation}
so that the world market absorbs any quantity at a price which is fixed in terms of foreign exchange.
Symmetrically, \textsc{row} supplies the varieties it exports to the modelled regions at a fixed price, $p^E_{is}=p^{FX}$ for $s\in R^X$, and these varieties enter the Armington import aggregates of the modelled regions as in Appendix \ref{equil}.
The market for foreign exchange clears through $p^{FX}$:
\begin{equation}\label{mktFX}
\sum_{r\in R^M}\textit{pem}^0_{ir}\textit{vem}_{ir}X_{ir} + \textit{rowpfx}
= \sum_{s\in R^X}\textit{vem}_{is}X_{is} + \sum_{j,s\in R^X}\textit{vst}_{js}\,\textit{YT}_j
+ \frac{\textit{ROW}}{p^{FX}},
\end{equation}
where the rest-of-world agent is endowed with an amount of foreign exchange, \textit{rowpfx} [\texttt{rowpfx}], provides the transport services that \textsc{row} regions export, and finances the current account balances of the modelled regions.
Its income is $\textit{ROW}=p^{FX}\textit{rowpfx} - p^I_n\sum_{r\in R^M}\textit{vb}_r - \sum_j p^T_j\sum_{i,s}\sum_{k\in R^X}\textit{vtwr}_{jisk}$ (the last term being the transport services used in shipments to \textsc{row} regions), which it spends on foreign exchange.
Equations \eqref{zpX}--\eqref{mktFX} are the closure of \citet{LanzRutherford2016}.
Without a finite elasticity of transformation between domestic and export supply ($\textit{etrndx}_i=\infty$, the \texttt{gtap} model version), \eqref{zpX} pins down the domestic producer price of every exported good, and output responds to trade policy in a bang-bang fashion.

\subsubsection{Small Armington trade economy, implementation \texttt{sate}}

In the \SATE closure, the export activity sells into a separate export market with price $p^{PX}_{ir}$,
\begin{equation}\label{zpXsate}
p^X_{ir}\,\frac{1-t^{xs,row}_{ir}}{1-\overline{t}^{xs,row}_{ir}} \;\geq\; p^{PX}_{ir}
\qquad \perp \; X_{ir}\geq 0,
\end{equation}
where the world's demand for the region's exports has the constant price elasticity $\sigma^m_i=\textit{esubm}_i$:
\begin{equation}\label{exportdemand}
E_{ir} = \left(\frac{p^{FX}}{p^{PX}_{ir}}\right)^{\sigma^m_i},
\qquad\quad
\textit{pem}^0_{ir}\textit{vem}_{ir}X_{ir} = \textit{pem}^0_{ir}\textit{vem}_{ir}E_{ir}
\quad \perp\; p^{PX}_{ir}.
\end{equation}
The representative agent of region $r$ sells the exports and receives their value in foreign exchange,
\begin{equation}\label{revexp}
\textit{RX}_{ir}\,p^{FX} = E_{ir}\,p^{PX}_{ir},
\end{equation}
so that the term $\sum_i \textit{pem}^0_{ir}\textit{vem}_{ir}\left(p^{FX}\textit{RX}_{ir}-p^{PX}_{ir}E_{ir}\right)=0$ is added to $\textit{RA}_r$ and the supply of foreign exchange in \eqref{mktFX} becomes $\sum_{r\in R^M}\textit{pem}^0_{ir}\textit{vem}_{ir}\textit{RX}_{ir}$.
In \MPSGE, $E_{ir}$ and $\textit{RX}_{ir}$ are auxiliary variables which ration an endowment of $-\textit{pem}^0_{ir}\textit{vem}_{ir}$ units of export goods and $+\textit{pem}^0_{ir}\textit{vem}_{ir}$ units of foreign exchange of agent \texttt{RA(r)} (see the \texttt{\$demand:RA(r)} blocks and the \texttt{\$constraint:EXPORTS} and \texttt{\$constraint:REVEXP} blocks in the code listing below).
Equations \eqref{zpXsate}--\eqref{revexp} require no information about the accounts of the rest of the world: the export demand curve is fully characterized by benchmark exports and the Armington elasticity.
Imports remain available at fixed world prices.

\subsubsection{Small Armington trade economy, implementation \texttt{sate2}}

Implementation \texttt{sate2} follows the large open economy formulation of \citet{Schreiber2024}.
The export activity combines the region's exports with a fixed factor owned by the rest of the world in a Cobb-Douglas technology and produces foreign exchange:
\begin{equation}\label{zpXsate2}
\left(p^X_{ir}\,\frac{1-t^{xs,row}_{ir}}{1-\overline{t}^{xs,row}_{ir}}\right)^{1-\theta^{x}_{i}}
\left(p^{FIX}_{ir}\right)^{\theta^{x}_{i}} \;\geq\; p^{FX}
\qquad \perp \; X_{ir}\geq 0,
\end{equation}
with the cost share of the fixed factor set to $\theta^x_i=1/\sigma^m_i$ [\texttt{theta\_xi(i)}] and a benchmark endowment of the fixed factor of $\textit{fix}^0_{ir}=\theta^x_i/(1-\theta^x_i)\,\textit{pem}^0_{ir}\textit{vem}_{ir}$ [\texttt{fix0(i,r)}], which accrues to the \textsc{row} agent.
With $p^{FX}$ given, the derived demand for the region's exports has price elasticity $1/\theta^x_i=\sigma^m_i$ \citep[Proposition 2]{Schreiber2024}, so that \texttt{sate} and \texttt{sate2} describe the same export demand curve.
The two implementations differ only in bookkeeping: in \texttt{sate2}, the rents associated with the downward sloping export demand curve are explicitly attributed to the rest of the world.
In the simulations reported in this paper, both implementations give identical results (up to the six digits reported).

\subsubsection{Average export taxes on exports to the rest of the world}

In the code of \citet{LanzRutherford2016}, $t^{xs,row}_{ir}$ is computed once by \texttt{loadrm} and does not follow subsequent changes of the bilateral rates $t^{xs}_{irs}$.
Program \texttt{tgrid.gms} now provides the macro \texttt{updrow(rr)} which recomputes $t^{xs,row}_{ir}$ after each change of export taxes (environment variable \texttt{--soextax=yes}, the default).
As a consequence, in the small open economy closures the counterfactual scaling of export taxes affects exports to the rest of the world, as it does in the \GMR model.
In the 2016 paper, only import tariffs moved in the \SOE simulations.
